%% title: Self-training: dose, labels against pseudo-labels, new traps
%% target: Sections/Results.tex
%% anchor: Images of species that already have a thousand buy accuracy on the saturated axis and cost it under shift.
%% action: after
%% why: The largest lever in the paper: unlabelled trap images, gated by confidence, gain 21 points on the probe set at no cost on the test fold and transfer to a new trap; real labels gain more on the labelled species and close the world on the rest; the dose optimum survives on unseen species only.
%% source: journal/2026-10-22-rescoring-in-the-mini-metrics-convention.md (probe table, amendment); data/rescore/probe-grouped.csv and probe-ho-grouped.csv (B1, B3rep1x/5x, B3, B3rep26x, lab500/2500/12230, T1rep, X1/X2 pooled, F1); paper/DRAFT.md §4.11 (error shares 6/9/26/59/91 %, teacher accuracy 98.15 %, in-distribution 0.9003)
%% check: Paste after step 19. | The interior optimum on the probe set, a headline of our earlier draft, does not survive the current convention; this piece states it for held-out species only. | Rerun: the dose sweep (1x, 11x, 26x) with the pinned recipe (decision sheet row 14). | The error-share percentages are counts of where errors land, independent of the convention.
%% status: ready
\subsection{Self-training on unlabelled trap images}\label{sec:results_selftraining}
\Cref{tab:selftraining} varies the target-domain data added to recipe L at 20\,M parameters.
Self-training with 12,230 gated pseudo-labels, each repeated 11 times (2.6\,\% of training rows), raises the probe set from 0.424 to 0.631 and its held-out species from 0.356 to 0.414, without a single trap label, and moves the test fold by 0.04 points (0.8999 to 0.9003).

\begin{table}[htb]
    \centering
    \caption{Target-domain data added to recipe L (EfficientNetV2-S with margin and trap-like augmentation, B1). Species macro-$\text{F}_1$, calibrated on 10\,\% of the (trap, night) groups.}
    \label{tab:selftraining}
    \begin{tabular}{lccc}
        \toprule
        trap data added & share of rows & probe & held-out species \\
        \midrule
        none & 0 & 0.424 & 0.356 \\
        12,230 pseudo-labels, once & 0.24\,\% & 0.531 & 0.389 \\
        12,230 pseudo-labels, 11 copies & 2.6\,\% & 0.631 & \textbf{0.414} \\
        12,230 pseudo-labels, 27 copies & 6.2\,\% & 0.624 & 0.348 \\
        12,230 pseudo-labels, 58 copies & 12.3\,\% & 0.606 & 0.306 \\
        \midrule
        500 / 2,500 real labels, once & & 0.488 / 0.492 & 0.357 / 0.342 \\
        12,230 real labels, once & 0.24\,\% & 0.609 & 0.379 \\
        12,230 real labels, 11 copies & 2.6\,\% & \textbf{0.723} & 0.312 \\
        \midrule
        pseudo-labels, 11 copies, trap held out & 2.6\,\% & 0.627 & 0.404 \\
        pseudo-labels, 11 copies, site held out & 2.6\,\% & 0.592 & 0.391 \\
        \bottomrule
    \end{tabular}
\end{table}

\emph{Dose.} One copy of each pseudo-labelled image, a quarter of a percent of the rows, buys half of the gain.
Beyond 11 copies the probe set stays flat (0.631, 0.624, 0.606) while the held-out species fall from 0.414 to 0.306: over-replication turns adaptation into memorisation of the pseudo-labelled species, visible only on the species that adaptation did not cover.

\emph{Labels against pseudo-labels.} At equal dose, real labels beat pseudo-labels (0.609 against 0.531 once; 0.723 against 0.631 with 11 copies).
That bounds the price of label noise from above, since the real labels also span every capture night while the gate keeps the teacher's easiest, 98\,\% correct, images.
But replicated real labels lose 10.2 points to self-training on the held-out species, falling below the model without any trap data.
The share of held-out species' errors that land on a species labelled at the traps rises with the label budget, from 6\,\% with no labels to 9, 26, 59 and 91\,\% for 500, 2,500, 12,230 and 12,230 labels with 11 copies; self-training sits at 42\,\%.
Replicated labels teach the model that anything on a trap sheet is one of the labelled species, which is the wrong lesson for a deployment that meets unknown species.

\emph{New traps.} Removing every pseudo-label of the trap a crop came from (twelve folds), or of its site's four traps (three folds), keeps 98\,\% and 81\,\% of the probe gain respectively.
Self-training adapts to capture conditions, not only to particular traps, and since it needs no labels it can be repeated for each new deployment.

\emph{Capacity.} The 20\,M model with self-training (0.631) is far above a 198\,M model without trap data (0.491) on the probe set, and slightly above it on held-out species (0.414 against 0.390).
Combined with the larger backbone and without oversampling (\cref{sec:results_longtail}), self-training gives B8 (0.647 and 0.452).
